%-------------------------------------------------------------------------------
%	SECTION TITLE
%-------------------------------------------------------------------------------
\cvsection{个人项目}

%-------------------------------------------------------------------------------
%	CONTENT
%-------------------------------------------------------------------------------
\begin{cventries}

%---------------------------------------------------------
  \cventry
    {AI项目} % 项目名称
    {模型训练、个人AI工具、开源贡献} % 角色
    {} % 地点
    {2023年至今} % 日期
    {
      \begin{cvitems} % 描述
        \item {使用nanoGPT在FineWeb数据集上从零训练GPT-2 124M两次——分别在RunPod H200和DigitalOcean H100上完成；另在配备RTX 4070的家用服务器上进行了实验。}
        \item {过去一年通过OpenRouter等LLM服务商消耗约15亿token；仅上月即消耗约5亿token。}
        \item {jekyll-ai-blog：AI驱动的博客平台，支持多语言自动翻译、Google Cloud TTS音频生成、XeLaTeX PDF/EPUB流水线及GitHub Actions工作流。}
        \item {lzwjava.github.io：个人博客与知识库，约400篇原创文章及约8000条AI回答笔记；过去一个月约7万页面浏览量（Cloudflare Analytics），访客最多的国家为新加坡。}
        \item {ww：跨平台CLI工具集——集成AI提交信息（Gemini Flash）、图片/PDF处理、网页搜索、GitHub Copilot Chat及LLM辅助功能。}
        \item {iclaw：终端AI智能体（REPL），可自主编码、搜索并执行Shell命令。支持GitHub Copilot（OAuth）和OpenRouter；适用于无IDE插件的受限企业环境。}
        \item {zz：ML项目数据集处理与训练工具集——涵盖FineWeb数据下载/提取、训练日志分析及GPT-2训练过程中使用的评估脚本。}
        \item {live-server现代化改造：通过AI编程工具完成Docker化，升级CodeIgniter和Vue至当前版本，实现全栈现代化。}
        \item {Tree\_Of\_Thought（1个PR）：为思维树推理系统贡献代码，新增OpenAI兼容请求模块及python-dotenv配置支持。}
      \end{cvitems}
    }

%---------------------------------------------------------
  \cventry
    {各类个人项目} % 项目名称
    {开源贡献、博客开发、后端与前端开发、移动应用开发、工具开发} % 角色
    {} % 地点
    {2013 - 2021} % 日期
    {
      \begin{cvitems} % 描述
        \item {算法解决方案：算法问题的解决方案，托管在GitHub上，有2466次提交，主要使用Java实现。（2013年至今）}
        \item {个人博客：维护一个双语博客（英文和中文），有500次提交，分享编程、技术和个人发展的见解。（2013年至今）}
        \item {直播服务器：使用PHP开发的知识直播平台后端，有660次提交。（2016 - 2017年）}
        \item {直播移动网页：使用Vue和JavaScript开发的知识直播平台移动前端，有528次提交。（2016 - 2017年）}
        \item {直播网页：使用Vue开发的知识直播平台桌面前端，有140次提交。（2016 - 2017年）}
        \item {直播微信小程序：使用JavaScript开发的知识直播平台微信小程序，有63次提交。（2016 - 2017年）}
        \item {代码审查服务器：使用PHP开发的专业代码审查平台后端，有275次提交。（2015 - 2016年）}
        \item {代码审查网页：使用Vue和JavaScript开发的专业代码审查平台前端，有302次提交。（2015 - 2016年）}
        \item {LeanChat iOS：展示LeanCloud SDK的iOS聊天应用，使用Objective-C实现，有556次提交。（2014 - 2015年）}
        \item {LeanChat Android：展示LeanCloud SDK的Android聊天应用，使用Java实现，有412次提交。（2014 - 2015年）}
        \item {费曼讲座Mobi：开发将LaTeX转换为SVG以创建mobi电子书的工具，使用Python实现，有47次提交。（2020 - 2021年）}
      \end{cvitems}
    }

%---------------------------------------------------------
\end{cventries}
